\documentclass[12pt,letterpaper]{article}


\usepackage[spanish]{babel}
\usepackage{caption}
\usepackage{csquotes} 
\usepackage[style=apa]{biblatex}
\usepackage{fontspec}
\usepackage{microtype}
\usepackage{float}
\usepackage{graphicx}
\usepackage{indentfirst}
\usepackage{setspace}
\usepackage{hyperref}
\usepackage[margin=2.54cm]{geometry}
\usepackage{lipsum}

\setmainfont{Arimo}
\doublespacing
\addbibresource{ref.bib}
\graphicspath{{images/}}

\hypersetup{
  colorlinks=true,
  linkcolor=black,
  citecolor=black,
  urlcolor=blue,
  pdfauthor={autor1, autor2, etc},
  pdftitle={Titulo},
  pdfsubject={Asignatura},
  pdfkeywords={Palabras clave}
}

\begin{document}

% ================= PORTADA =================
\thispagestyle{empty}

\begin{center}
    {\LARGE\textbf{TÍTULO DEL PROYECTO}}\\[2cm]

    \textbf{Estudiantes:}\\[0.5cm]
    NOMBRES 1 \\
    NOMBRES 2 \\[3cm]

    \includegraphics[height=4cm]{logo.png}\\[3cm]

    \textbf{Universidad del Quindío} \\
    \textbf{Facultad Ingeniería} \\
    \textbf{Ingeniería de Sistemas y Computación} \\
    NOMBRE DE LA ASIGNATURA \\
    Armenia, Quindío\\
    MES Y AÑO
\end{center}

\newpage

% ================= CUERPO =================
\pagestyle{plain}

\tableofcontents
\newpage

% ================= INTRODUCCIÓN =================
\addcontentsline{toc}{section}{Introducción}

\lipsum[1] % REEMPLAZAR con texto real

No olvidar usar referencias \parencite{name}

% ================= SECCIÓN 1 =================
\section{NOMBRE DE LA PRIMERA SECCIÓN}

\lipsum[2]

\subsection{NOMBRE DE LA SUBSECCIÓN}
\lipsum[3]

% ================= SECCIÓN 2 =================
\section{NOMBRE DE LA SEGUNDA SECCIÓN}

\lipsum[4]

\subsection{NOMBRE DE LA SUBSECCIÓN}
\lipsum[5]


% ================= CONCLUSIONES =================
\section{Conclusiones}

\lipsum[8]

% ================= BIBLIOGRAFÍA =================
\newpage
\printbibliography

\end{document}
